\chapter{绪论}
\label{chap:intro}

\section{研究背景}

阅读正在从纸面迁移到屏幕。新闻、小说、技术文档乃至学术论文，
越来越多地以长文本的形式出现在手机、平板和显示器上。
与纸书不同，屏幕上的版式由软件决定：行长随窗口变化，行距由样式表设定，
同一篇文章在不同设备上可能呈现出完全不同的面貌。

可读性研究已有近百年历史\cite{flesch1948}，但经典公式主要面向英文短文，
并且默认版面是固定的。中文没有词间空格，长文本又常常跨越数十屏，
读者的疲劳更多来自「找下一行」「找回刚才读到的位置」，而这些恰恰是版面造成的。

\section{研究问题}

本文试图回答三个问题：
\begin{enumerate}
  \item 在中文长文本中，版面因素对阅读负荷的影响有多大？
  \item 能否用少数几个可计算的指标刻画这种影响？
  \item 把版面指标加入可读性模型后，评估结果能提升多少？
\end{enumerate}

\section{主要贡献}

\begin{enumerate}
  \item 构建了中文长文本阅读数据集，包含 1\,200 篇文章、312 位读者、36\,000 条阅读记录；
  \item 提出版面负荷指标 $\Lambda$，由行长、行距与段落密度三部分组成；
  \item 设计了融合语言与版面特征的评估模型，并开源全部数据与代码。
\end{enumerate}

\section{论文结构}

第~\ref{chap:related}~章回顾相关研究；第~\ref{chap:method}~章介绍版面负荷指标与评估模型；
第~\ref{chap:experiment}~章报告实验结果；第~\ref{chap:conclusion}~章总结全文并讨论后续工作。
